\begin{solapartat}
\[ {}^{240}_{94}\,\mathrm{Pu} + {}^{a}_{b}\,\mathrm{x} \rightarrow {}^{241}_{94}\,\mathrm{Pu} \Rightarrow
\left\{\begin{aligned} 240 + a &= 241 \Rightarrow a = 1 \ \text{\punts{0,1}}\\ 94 + b &= 94 \Rightarrow b = 0 \ \text{\punts{0,1}}\end{aligned}\right. \]
\[ {}^{241}_{94}\,\mathrm{Pu} \rightarrow {}^{241}_{95}\,\mathrm{Am} + {}^{c}_{d}\,\mathrm{y} \Rightarrow
\left\{\begin{aligned} 241 &= c + 241 \Rightarrow c = 0 \ \text{\punts{0,1}}\\ 94 &= d + 95 \Rightarrow d = -1 \ \text{\punts{0,1}}\end{aligned}\right. \]

El ${}^{240}$Pu ha capturat un neutró \punts{0,3}

El ${}^{241}$Pu ha emes una partícula $\beta$ ó electró. La desintegració s'anomena emissió beta \punts{0,3}
\end{solapartat}

\begin{solapartat}
La llei de desintegració la podem escriure com:
\[ N(t) = N_0\,\mathrm{e}^{-\frac{t\,\ln 2}{t_{1/2}}} = N_0\,\left(\mathrm{e}^{-\ln 2}\right)^{\frac{t}{t_{1/2}}} = N_0\,2^{-\frac{t}{t_{1/2}}} \quad \text{\punts{0,5}} \]
\% de nuclis que s'hauran desintegrat després de (2013-1944) = 69 anys:
\[ 100\cdot\frac{N_0 - N(t = 69)}{N_0} = 100\cdot\frac{N_0\,\left(1 - 2^{-\frac{69}{432}}\right)}{N_0} = 10,5\,\% \quad \text{\punts{0,5}} \]
\end{solapartat}
